% !PRINTED 268
\weberpage{268}
%<<CONT
\[
\alpha^2 - \gamma = \delta^2,
\]
so ist $\alpha^2 > \delta^2$, also, wenn $\delta$ positiv genommen wird
\[
\alpha > \delta
\]
und
\[
27 D = 4 \delta^6 - (2\alpha\delta^2 - \beta)^2.
\]
Hiernach kann $D$ bei negativem $\beta$ nicht positiv sein; denn ist $\beta$ negativ, so ist
\[
2\alpha\delta^2 - \beta > 2\alpha\delta^2 > 2\delta^3,
\]
also $4\delta^6 - (2\alpha\delta^2 - \beta)^2$ negativ.

Es ist hiermit direct nachgewiesen, dass die Kriterien (5) und (6) genau dasselbe besagen.

Nach der im \S.\,78 gegebenen geometrischen Interpretation bedeutet $\beta = 0$ eine Fläche dritter Ordnung, die durch die Parabel $b = 0$, $\gamma = 0$ hindurchgeht, und die, soweit negative Werthe von $a$ in Betracht kommen, ganz in dem Raumtheil verläuft, in dem $D$ negativ ist.

Um für die Gleichung fünften Grades ein Beispiel vor Augen zu haben, betrachten wir die Gleichung
\[
x^5 + x^2 + a = 0.
\]
Wir haben also in unseren allgemeinen Ausdrücken \S.\,73 oder auch \S.\,74, (2), (3) zu setzen
\[
a_0, a_1, a_2, a_3, a_4, a_5 =
\]
\[
1, 0, 0, 1, 0, a,
\]
und wir erhalten für die Determinante der Bezoutiante
\[
\begin{vmatrix}
-2a, & 0, & 0, & -5a \\
0, & \tfrac{6}{5}, & -5a, & 0 \\
0, & -5a, & 0, & -3 \\
-5a, & 0, & -3, & 0
\end{vmatrix}
\]
und für die Discriminante
\[
D = 108 a + 3125 a^4.
\]
Eine Kette von Haupt-Unterdeterminanten ist
\weq{8}{D, \; 50a^3, \; \frac{-12a}{5}, \; -2a, \; 1.}
Die Discriminante ist negativ, wenn
\weq{9}{0 < -a < \sqrt[3]{\frac{108}{3125}},}
sonst positiv, besonders also für alle positiven $a$ positiv. Man sieht, dass bei positiver Discriminante in (8) immer zwei
%CONT>>
